\section{Gaussian Space --- Invariance --- Coordinates}
%%% generated by the blueprint pipeline from TCSlib.BooleanAnalysis.GaussianSpace.Invariance.Coordinates
\subsection*{Gaussian Space --- Invariance --- Coordinates: Introduction}
A multilinear polynomial splits into a part omitting a chosen coordinate and that coordinate times its derivative. Coefficient, influence, degree, and independence identities support coordinate replacement.

\subsection*{Gaussian Space --- Invariance --- Coordinates: Results}

\begin{definition}[Polynomial part omitting a coordinate]
\lean{BooleanAnalysis.GaussianSpace.Invariance.coordinateFree}
\leanok
\uses{BooleanAnalysis.ThresholdFunctions.MultilinearPolynomial}
Let $F(x)=\sum_{S\subseteq\{1,\ldots,n\}}a_S\prod_{j\in S}x_j$ be a real multilinear polynomial and fix $i\in\{1,\ldots,n\}$. The coordinate-free part at $i$ is the multilinear polynomial with coefficient $0$ on sets containing $i$ and coefficient $a_S$ on every set $S$ not containing $i$.
\end{definition}

\begin{definition}[Coefficient polynomial of a cube derivative]
\lean{BooleanAnalysis.GaussianSpace.Invariance.coordinateDerivative}
\leanok
\uses{BooleanAnalysis.ThresholdFunctions.MultilinearPolynomial, BooleanAnalysis.ThresholdFunctions.MultilinearPolynomial.eval, BooleanAnalysis.derivative, BooleanAnalysis.fourierCoeff}
Let $F(x)=\sum_{S\subseteq\{1,\ldots,n\}}a_S\prod_{j\in S}x_j$ be a real multilinear polynomial and fix $i\in\{1,\ldots,n\}$. The coordinate derivative polynomial has, at each set $S$, the Fourier coefficient indexed by $S$ of the Boolean-cube derivative of $F$ restricted to sign inputs.
\end{definition}

\begin{lemma}[Coefficients of a coordinate derivative]
\lean{BooleanAnalysis.GaussianSpace.Invariance.coordinateDerivative_coeff}
\leanok
\uses{BooleanAnalysis.GaussianSpace.Invariance.coordinateDerivative, BooleanAnalysis.Hypercontractivity.fourierCoeff_derivative, BooleanAnalysis.ThresholdFunctions.MultilinearPolynomial, BooleanAnalysis.ThresholdFunctions.MultilinearPolynomial.fourierCoeff_eval}
\difficulty{1}
Let $F$ be a real coefficient family on the subsets of $[n]=\{0,\ldots,n-1\}$, and let
$i\in[n]$. The coefficient at $S\subseteq[n]$ of the polynomial formed from the Fourier
coefficients of the Boolean-cube derivative of $F$ is $0$ if $i\in S$, and is
$F(S\cup\{i\})$ if $i\notin S$.
\end{lemma}

\begin{lemma}[Degree reduction under coordinate differentiation]
\lean{BooleanAnalysis.GaussianSpace.Invariance.coordinateDerivative_degree}
\leanok
\uses{BooleanAnalysis.GaussianSpace.Invariance.coordinateDerivative, BooleanAnalysis.GaussianSpace.Invariance.coordinateDerivative_coeff, BooleanAnalysis.ThresholdFunctions.MultilinearPolynomial, BooleanAnalysis.ThresholdFunctions.MultilinearPolynomial.HasDegreeAtMost}
\difficulty{4}
Let $F$ be a real coefficient family on subsets of $[n]$, let $i\in[n]$, and let
$k\in\mathbb N$. If $F(S)=0$ whenever $|S|>k+1$, then the coefficient family formed from
the Fourier coefficients of the Boolean-cube derivative of $F$ has degree at most $k$;
that is, its coefficient at $S$ vanishes whenever $|S|>k$.
\end{lemma}

\begin{lemma}[Evaluation of the derivative polynomial]
\lean{BooleanAnalysis.GaussianSpace.Invariance.coordinateDerivative_eval}
\leanok
\uses{BooleanAnalysis.GaussianSpace.Invariance.coordinateDerivative, BooleanAnalysis.ThresholdFunctions.MultilinearPolynomial, BooleanAnalysis.ThresholdFunctions.MultilinearPolynomial.eval, BooleanAnalysis.derivative, BooleanAnalysis.walsh_expansion}
\difficulty{3}
Let $F$ be a real coefficient family on subsets of $[n]$ and let $i\in[n]$. At every
Boolean-cube input, the Walsh-basis evaluation of the coefficient family formed from the
Fourier coefficients of the Boolean-cube derivative of $F$ equals that discrete
derivative.
\end{lemma}

\begin{lemma}[Derivative coefficient energy and influence]
\lean{BooleanAnalysis.GaussianSpace.Invariance.coordinateDerivative_sum_sq}
\leanok
\uses{BooleanAnalysis.GaussianSpace.Invariance.coordinateDerivative, BooleanAnalysis.GaussianSpace.Invariance.coordinateDerivative_eval, BooleanAnalysis.ThresholdFunctions.MultilinearPolynomial, BooleanAnalysis.ThresholdFunctions.MultilinearPolynomial.eval, BooleanAnalysis.ThresholdFunctions.MultilinearPolynomial.expect_eval_sq, BooleanAnalysis.ThresholdFunctions.influence_eq_expect_derivative_sq, BooleanAnalysis.influence}
\difficulty{3}
Let $F$ be a real coefficient family on subsets of $[n]$, and let $i\in[n]$. The sum,
over all $S\subseteq[n]$, of the squares of the coefficients formed from the Fourier
coefficients of the Boolean-cube derivative of $F$ equals the influence of coordinate
$i$ for the Boolean-cube function obtained by evaluating $F$ in the Walsh basis.
\end{lemma}

\begin{lemma}[Evaluation of the coordinate-free polynomial]
\lean{BooleanAnalysis.GaussianSpace.Invariance.coordinateFree_eval_sum}
\leanok
\uses{BooleanAnalysis.GaussianSpace.Invariance.coordinateFree, BooleanAnalysis.Hypercontractivity.MultilinearMoments.eval_supported, BooleanAnalysis.Hypercontractivity.evalMultilinearRV, BooleanAnalysis.ThresholdFunctions.MultilinearPolynomial}
\difficulty{4}
Let $F$ be a real coefficient family on subsets of $[n]$, let $i\in[n]$, let $\Omega$ be
a measurable space, let $X_j:\Omega\to\mathbb R$ for each $j\in[n]$, and let
$\omega\in\Omega$. The random-input evaluation of the polynomial obtained from $F$ by
setting to zero every coefficient indexed by a set containing $i$ equals
$\displaystyle\sum_{T\subseteq[n]\setminus\{i\}}F(T)\prod_{j\in T}X_j(\omega)$.
\end{lemma}

\begin{lemma}[Evaluation of the derivative on random inputs]
\lean{BooleanAnalysis.GaussianSpace.Invariance.coordinateDerivative_eval_sum}
\leanok
\uses{BooleanAnalysis.GaussianSpace.Invariance.coordinateDerivative, BooleanAnalysis.GaussianSpace.Invariance.coordinateDerivative_coeff, BooleanAnalysis.Hypercontractivity.MultilinearMoments.eval_supported, BooleanAnalysis.Hypercontractivity.evalMultilinearRV, BooleanAnalysis.ThresholdFunctions.MultilinearPolynomial}
\difficulty{4}
Let $F$ be a real coefficient family on subsets of $[n]$, let $i\in[n]$, let $\Omega$ be
a type, let $X_j:\Omega\to\mathbb R$ for each $j\in[n]$, and let $\omega\in\Omega$. The
multilinear evaluation of the coefficient family formed from the Fourier coefficients of
the Boolean-cube derivative is $\displaystyle\sum_{T\subseteq[n]\setminus\{i\}}
F(T\cup\{i\})\prod_{j\in T}X_j(\omega)$.
\end{lemma}

\begin{lemma}[Coordinate decomposition of polynomial evaluation]
\lean{BooleanAnalysis.GaussianSpace.Invariance.eval_coordinate_split}
\leanok
\uses{BooleanAnalysis.GaussianSpace.Invariance.coordinateDerivative, BooleanAnalysis.GaussianSpace.Invariance.coordinateDerivative_eval_sum, BooleanAnalysis.GaussianSpace.Invariance.coordinateFree, BooleanAnalysis.GaussianSpace.Invariance.coordinateFree_eval_sum, BooleanAnalysis.Hypercontractivity.evalMultilinearRV, BooleanAnalysis.Hypercontractivity.multilinear_sum_insert, BooleanAnalysis.ThresholdFunctions.MultilinearPolynomial}
\difficulty{3}
Let $n$ be a natural number, let $F$ be a real multilinear polynomial on the Boolean
cube of dimension $n$, let $i$ be a coordinate, and let $X_j:\Omega\to\mathbb R$ be a
family of real-valued functions on a type of inputs $\Omega$. For every
$\omega\in\Omega$, the evaluation of $F$ at $X(\omega)$ is the sum of its part omitting
coordinate $i$ and $X_i(\omega)$ times the evaluation of the coefficient polynomial of
the cube derivative in coordinate $i$.
\end{lemma}

\begin{lemma}[Independence from coordinate-free polynomial parts]
\lean{BooleanAnalysis.GaussianSpace.Invariance.indep_coordinate_parts}
\leanok
\uses{BooleanAnalysis.GaussianSpace.Invariance.coordinateDerivative, BooleanAnalysis.GaussianSpace.Invariance.coordinateDerivative_eval_sum, BooleanAnalysis.GaussianSpace.Invariance.coordinateFree, BooleanAnalysis.GaussianSpace.Invariance.coordinateFree_eval_sum, BooleanAnalysis.Hypercontractivity.MultilinearMoments.indep_sum_of_notMem, BooleanAnalysis.Hypercontractivity.evalMultilinearRV, BooleanAnalysis.ThresholdFunctions.MultilinearPolynomial}
\difficulty{3}
Let $n$ be a nonnegative integer, let $(\Omega,\mu)$ be a measured space, and let
$X_1,\ldots,X_n$ be an independent family of real-valued functions, each measurable up
to a null set. For a real coefficient family $F$ on subsets of the coordinates and a
fixed coordinate $i$, form its coordinate-free part and its coordinate-derivative
coefficient family as defined above. Then $X_i$ is independent of the pair consisting of
the multilinear evaluations of those two coefficient families at $X$.
\end{lemma}

%%% added by the blueprint pipeline
\subsection*{Gaussian Space --- Invariance --- Coordinates: Further results}

\begin{lemma}[Coordinate independence of polynomial parts]
\lean{BooleanAnalysis.GaussianSpace.Invariance.coordinate_parts_eq_of_off_coordinate}
\label{BooleanAnalysis.GaussianSpace.Invariance.coordinate_parts_eq_of_off_coordinate}
\leanok
\uses{BooleanAnalysis.GaussianSpace.Invariance.coordinateDerivative, BooleanAnalysis.GaussianSpace.Invariance.coordinateDerivative_coeff, BooleanAnalysis.GaussianSpace.Invariance.coordinateFree, BooleanAnalysis.Hypercontractivity.MultilinearMoments.eval_eq_of_off_coordinate, BooleanAnalysis.Hypercontractivity.evalMultilinearRV, BooleanAnalysis.ThresholdFunctions.MultilinearPolynomial}
\difficulty{3}
Let $F$ be a real coefficient family on subsets of $[n]$, let $i\in[n]$, let $\Omega$ be
a type, and let $X,Y:[n]\to(\Omega\to\mathbb R)$. Suppose $X_j=Y_j$ as functions on
$\Omega$ for every $j\ne i$. Then, as functions on $\Omega$, the random-input
evaluations of the coordinate-free part of $F$ agree for $X$ and $Y$, and the
random-input evaluations of the coefficient family formed from the Fourier coefficients
of the Boolean-cube derivative of $F$ also agree for $X$ and $Y$.
\end{lemma}
